\section{By Example: Implementing the Proposal}\label{sec:by_example_implementing_the_proposal}

% --------------------- Sub Section 1 ---------------------

\subsection[ ]{Motivational Example}\label{subsec:example_motivational_example}

\begin{frame}{\secname}
	\framesubtitle{\subsecname}

	To illustrate how the proposed approach works in practice, let's revisit the motivational example we presented earlier.

	\begin{figure}[h!]
		\sliderust*[0.98\linewidth]{motivational-example-solution.rs}
	\end{figure}
	% As we can see, the changes are minimal and mostly involve adding the necessary annotations to support the new functionality.
	% This demonstrates the practicality and ease of adoption of the proposed solution.
	% The spec! macro requires to specify all the types of the items involved in the method call, but we'll discuss why this is necessary later.
	% The #when macro is used to specify the conditions under which this implementation is valid.

\end{frame}

% --------------------- Sub Section 2 ---------------------

\subsection[ ]{Supported Predicates and Forms}\label{subsec:example_supported_predicates_and_forms}

\begin{frame}{\secname}
	\framesubtitle{\subsecname}

	The following predicates are supported:

	\begin{tabular}{p{0.45\textwidth} p{0.55\textwidth}}
		\begin{itemize}
			\item \inlinerust{T = Type} % Type equality: a generic type parameter is of a specific type
			\item \inlinerust{all(p1, p2, ...)} % Logical AND: all inner predicates are satisfied
			\item \inlinerust{not(p)} % Logical NOT: the inner predicate is not satisfied
		\end{itemize}
		 &
		\begin{itemize}
			\item \inlinerust{T: Tr1 + Tr2} % Trait implementation: a generic type parameter implements a specific trait or multiple traits
			\item \inlinerust{any(p1, p2, ...)} % Logical OR: at least one of the inner predicates is satisfied
		\end{itemize}
	\end{tabular}
	% These predicates can be combined to form complex conditions that must be met for the implementation to be valid.

	% \pause
	\bigskip

	We also support the following implementations forms:

	\begin{itemize}
		\item \inlinerust{impl<T> Trait for T} % Generic impl type: the type being implemented is a generic type parameter
		\item \inlinerust{impl<T> Trait for Type<T>} % Generic type parameter: the type being implemented is a concrete type with a generic type parameter
		\item \inlinerust{impl<T> Trait<T> for Type} % Generic trait parameter: the trait being implemented has a generic type parameter
	\end{itemize}
	% And all combinations of the above forms.

\end{frame}

% --------------------- Sub Section 3 ---------------------

\subsection[ ]{Lifetimes}\label{subsec:example_lifetimes}

\begin{frame}{\secname}
	\framesubtitle{\subsecname}
	\lipsum[1-1]
\end{frame}

% --------------------- Sub Section 4 ---------------------

\subsection[ ]{}\label{}

\begin{frame}{\secname}
	\framesubtitle{\subsecname}
	\lipsum[1-1]
\end{frame}
